\section{Machine Learning Domain and Methodological Foundations}
\label{sec:ml_methodology}

The core analytical capability of HormuzWatch is driven by a specialized unsupervised anomaly detection and probability calibration engine. This section outlines the theoretical, mathematical, and operational justifications for our methodological choices.

\subsection{The Unsupervised Learning Paradigm in Non-Cooperative Littorals}
A primary design question in military and maritime AI systems is whether to implement supervised classification (e.g., training a deep neural network or gradient-boosted decision tree on labeled attack incidents) or unsupervised anomaly detection.

HormuzWatch deliberately adopts an **unsupervised paradigm** based on three operational realities:
\begin{enumerate}[leftmargin=*]
    \item \textbf{Extreme Peacetime Class Imbalance:} Genuine maritime security incidents (e.g., vessel boardings, mine laying, naval interceptions, weapon smuggling) constitute less than $0.001\%$ of all broadcast AIS pings. Supervised models trained on such severely imbalanced data suffer from severe overfitting and collapse into majority-class classifiers.
    \item \textbf{Non-Stationary Adversarial Tactics:} Hostile actors dynamically modify their tactical profiles to evade static detection signatures. A supervised classifier trained on 2019 Gulf of Oman tanker limpet mine attacks will fail to generalize to novel uncrewed surface vessel (USV) harassment or GPS spoofing tactics deployed in 2026.
    \item \textbf{High-Volume Normative Baselines:} Conversely, normative commercial navigation is abundant. Tens of thousands of routine transits through the Strait of Hormuz define a dense, predictable manifold of physical kinematic trajectories. Deviations from this normative manifold can be robustly quantified without requiring historical attack labels.
\end{enumerate}

\subsection{Dual-Algorithm Ensembling: Isolation Forest and Local Outlier Factor}
Rather than relying on a single anomaly detection algorithm, HormuzWatch implements a complementary ensemble combining **Isolation Forest (iForest)** with **Local Outlier Factor (LOF)**.

\subsubsection{Isolation Forest: Global Hyperplane Isolation}
Isolation Forest isolates anomalies by recursively generating random axis-aligned partitioning hyperplanes in the $d$-dimensional feature space. Because anomalous observations are few and physically distant from normative clusters, they are isolated close to the root of the tree:
\begin{equation}
s_{\text{IF}}(x, n) = 2^{-\frac{E(h(x))}{c(n)}}
\end{equation}
where $h(x)$ is the path length of observation $x$ in a tree of $n$ instances, $E(h(x))$ is the expectation across an ensemble of 200 random isolation trees, and $c(n)$ is the average path length of an unsuccessful search in a Binary Search Tree:
\begin{equation}
c(n) = 2 \ln(n - 1) + 0.5772156649 \text{ (Euler-Mascheroni constant)} - \frac{2(n - 1)}{n}
\end{equation}
\textit{Algorithmic Advantage:} Global search efficiency with $O(n \log n)$ fitting complexity and rapid sub-millisecond inference.

\textit{Algorithm Blind Spot:} Vulnerable to axis-aligned masking effects and incapable of detecting localized density variations within non-convex clusters.

\subsubsection{Local Outlier Factor: Neighborhood Density Estimation}
Local Outlier Factor measures the local density of an observation with respect to its $k$-nearest neighbors ($k=20$). For a given point $p$:
\begin{enumerate}[leftmargin=*]
    \item The $k$-distance of $p$, denoted $d_k(p)$, is the Euclidean distance $d(p, o)$ to its $k$-th nearest neighbor.
    \item The reachability distance of point $p$ with respect to neighbor $o$ is defined as:
    \begin{equation}
    \text{reach-dist}_k(p, o) = \max\{d_k(o), d(p, o)\}
    \end{equation}
    \item The local reachability density ($\text{lrd}$) of $p$ is the inverse of the average reachability distance over its $k$-nearest neighbors $N_k(p)$:
    \begin{equation}
    \text{lrd}_k(p) = \left[ \frac{\sum_{o \in N_k(p)} \text{reach-dist}_k(p, o)}{|N_k(p)|} \right]^{-1}
    \end{equation}
    \item The Local Outlier Factor ($\text{LOF}$) compares the $\text{lrd}$ of $p$ against the $\text{lrd}$ of its neighbors:
    \begin{equation}
    \text{LOF}_k(p) = \frac{\sum_{o \in N_k(p)} \frac{\text{lrd}_k(o)}{\text{lrd}_k(p)}}{|N_k(p)|}
    \end{equation}
\end{enumerate}
\textit{Algorithmic Advantage:} Accurately detects local structural outliers—such as a cargo vessel moving at nominal speed but heading in reverse within an active Traffic Separation Scheme corridor.

\subsubsection{Ensemble Blending Formulation}
Raw scores from both algorithms are projected onto $[0, 1]$ using empirical bounds derived exclusively from the training set, followed by linear blending:
\begin{equation}
s_{\text{ens}}(x) = 0.55 \cdot \text{norm}(s_{\text{IF}}(x)) + 0.45 \cdot \text{norm}(s_{\text{LOF}}(x))
\end{equation}
The $0.55 / 0.45$ weighting provides a balance: Isolation Forest anchors global kinematic boundaries (e.g., impossible speeds or wild course jumps), while LOF sensitizes the system to micro-spatial deviations within dense maritime corridors.

\subsection{Probability Calibration: Isotonic Regression vs. Platt Scaling}
Raw anomaly scores $s_{\text{ens}} \in [0, 1]$ represent relative rank-order outlierness rather than true probabilities. Deploying raw scores directly into an operations center produces extreme miscalibration.

\subsubsection{Why Platt Scaling (Logistic Sigmoid) Was Rejected}
Platt scaling fits a parametric logistic regression to score outputs:
\begin{equation}
P(Y = 1 \mid s) = \frac{1}{1 + \exp(A \cdot s + B)}
\end{equation}
This formulation enforces a strict, symmetric sigmoidal curve. In maritime anomaly detection, the boundary between normative vessel clusters and outlying behavior is sharp, asymmetric, and multi-modal. Forcing a rigid sigmoid curve onto complex multi-density distributions compresses tail probabilities and distorts intermediate risks.

\subsubsection{The Isotonic Regression Advantage}
HormuzWatch employs non-parametric **Isotonic Regression**, which fits a free-form, non-decreasing step function using the Pool Adjacent Violators Algorithm (PAVA):
\begin{equation}
\min_{\hat{m}} \sum_{i=1}^{N_{\text{calib}}} \left( y_i - \hat{m}(s_{\text{ens}, i}) \right)^2 \quad \text{subject to } \hat{m}(u) \le \hat{m}(v) \text{ whenever } u \le v
\end{equation}
\textit{Operational Benefits:}
\begin{itemize}[leftmargin=*]
    \item \textbf{Distribution-Agnostic:} Adapts directly to the empirical error distribution without parametric distortion.
    \item \textbf{Guaranteed Monotonicity:} A vessel with a higher combined outlier score $s_A > s_B$ is mathematically guaranteed to receive a calibrated probability $P_A \ge P_B$.
    \item \textbf{Empirical Superiority:} Reduced Expected Calibration Error (ECE) from $14.86\%$ down to $4.57\%$, representing a \textbf{69.2\% improvement} in probability reliability.
\end{itemize}

\subsection{Operational Distinction: Anomaly Probability vs. Rogue Behavior}
A foundational principle enforced across the HormuzWatch codebase is the technical and legal separation between kinematic outlier probability and hostile intent:
\begin{equation}
P(\text{Kinematic Outlier} \mid \text{Telemetry}) \neq P(\text{Hostile / Rogue Intent} \mid \text{Kinematic Outlier})
\end{equation}
A merchant tanker executing an abrupt $90^\circ$ heading change at $20\text{ kts}$ is genuinely a statistical outlier ($P(\text{Anomaly}) \approx 98\%$). However, under maritime law (COLREGs Rules 8 and 17), such an action may represent a legal emergency maneuver to avoid collision with a submerged reef or unlit fishing skiff. Labeling such an event as ``$98\%$ Rogue Behavior'' would corrupt operational reporting.

HormuzWatch strictly labels this signal as **Calibrated Anomaly Probability**, which is subsequently synthesized within the multi-factor Intelligence Fusion Engine alongside geopolitical threat intelligence and deterministic rule checks.
